% New proof, 2026-10-22. No historical-byte identity is claimed.
\section{R2 with an unbounded source ratio and a bounded relative current}
\label{sec:r2-thin-new}

\paragraph{Statement and domain.}
Use the actual original weak-C2 source RER1--RER8, normalized by
$d=1$, with $\lambda=\beta=0$ and every original strict gap.
Suppose $t,a$ are positive unbounded, while $B,aB,tB$ are
positive infinitesimals, and put $e=1-r$.  Assume
\[
             \lambda_n=n/(tB)=O(1),\qquad n\geq0.       \tag{TU1}
\]
The entire original parent is tensor SOS for every independent
projective test and every state.  Literal $n=0$ is included.
No comparison of a testing quantity with a source quantity is made.
This theorem treats precisely this source sector; it is not a
completion of the other R2 sectors or the global finite-lift proof.

The bounded-Gram and finite toric allocation lemmas used below are
the actual retained SC16--SC20 and QI1--QI10 proofs.  We use their
printed finite algebra and prove its source bounds on the present
domain.  We do not apply the physical C1 theorem to a formal source
whose event coefficient could be negative.

\subsection{An exact common contact and bounded source comparison}

Set
\[
 S=1/t,\quad z=1/a,\quad q=a(tB),\quad w=ae,\qquad
 b=\sqrt S,\quad g=\sqrt z,\quad c=\sqrt q.
                                                               \tag{TU2}
\]
Thus $B=Szq$, $e=zw$, the actual squared angles are
$u_c^2=Sq$ and $v_c^2=zq$, and $S,z,q$ are positive
infinitesimals.  Keep $h$ as the original C2 source parameter.
The strict $m$ gap gives
\[
 R=1+Sqh,\quad M=R-w>0,
 \qquad 0<w<R\sim1.                                      \tag{TU3}
\]
Both original openings $h,1-h$ remain arbitrary positive quantities.

The exact RER2 identity, divided by $e>0$, gives
\[
 Q_{\rm RER}
 =\tau\{(1+z)(1+Sh)(1+Sq)+2r(z\Delta)\lambda_n\},
 \qquad \tau=q/(zw),\quad \lambda_n=n/(qz).
\]
Here $Q_{\rm RER}\sim1$, $z\Delta\sim1$ and $r\sim1$.
The bracket is a positive bounded unit, since $\lambda_n\geq0$
is bounded and its first term is a positive unit.  Consequently
\[
 \tau\text{ is a positive bounded unit},\qquad
 \overline\tau=\frac1{1+2\overline\lambda_n},\qquad
 q=zw\tau.
                                                               \tag{TU4}
\]
No division by $n$ has occurred.  This sector includes the earlier
thin condition $\rho_n=n/e^2\in\mathrm I$: in that case the same
RER identity gives $\tau\sim1$ and
$\lambda_n=\rho_nw/\tau\in\mathrm I$ directly, without first
assuming (TU1).

For the comparison only, substitute into the native C1 rational
formulas NC1--NC4
\[
 k^2=Sq,\quad \eta=1/S,\quad \zeta=z,\quad
 r_{\rm C1}=h,\quad t_{\rm C1}=w,\quad
 a_{\rm C1}=S(1-h).                                     \tag{TU5}
\]
Their $\theta$ is $e=zw$ and their $R,M$ are exactly (TU3).
Write a superscript $0$ for these formal C1 coefficients.
All their source denominators have the requisite positive signs:
\[
\begin{gathered}
 K=1+Sq,\quad J=1+Sq(1-e)-S(1-h),\quad
 E=1-q(1-e),\quad Q=1-Sqzh,
\end{gathered}
\]
and $J,E,Q$ are positive bounded units.  Their common polynomial is
\[
 P=1+qz+Sq(2+z-2e)+Sq^2(z+e-e^2)+S^2hq^2(1+z-e).
                                                               \tag{TU6}
\]
In particular $P\sim1$.  We make no sign claim for $\beta^0$.

The two sources have exactly the same state contact.  With $Z_c=1$,
\[
 u_c=bc,\quad v_c=cg,\quad W_c=bg,\quad
 Y_c=hu_cW_c,\quad X_c=-(1-e)u_c.                       \tag{TU7}
\]
This follows either by substitution in the two native dictionaries,
or from $Y_c=hu_cW_c$, $X_c=-ru_cZ_c$ and
$Z_c/W_c=\sqrt{at}$ in the actual C2 source.  In particular
the comparison does not move the contact.

Use the actual, not formal, normalized coefficients
\[
 c_b=SC,\ a_b=SA,\ D_b=D/S,\ \alpha_b=\alpha/(Sz),\quad
 h_2=mzM,\quad h_b=zw/2.
                                                               \tag{TU8}
\]
They satisfy $m=h_2/(zM)$ and $\beta+1/2=h_b/(zw)$.
Use the same normalizations with superscript $0$, and let $\delta$
mean actual minus formal.  The following explicit polynomials make
all needed cancellations finite and opening-safe.  Put $G=1+z-e$,
and define
\begin{align*}
 P_C={}&1+qz+S\{1+q(2+2z-e)+q^2z\}\\
 &+S^2\{q[h(1+z)+1-e]
          +q^2[hez-e(1+z)+2z+1]\}+S^3hq^2G,\\
 P_A={}&1+qz+S\{q(2+z-2e)+q^2z(1-e)\}
                 +S^2q^2G(G-hz),\\
 P_D={}&1+q(z+e)+S\{q[h(1+z)+1-e]+q^2(z+e-e^2)\}
                 +S^2hq^2G,\\
 P_m={}&1+qz+S\{q(2+z-2e)+q^2z\}
                 +S^2q^2\{(1-e)^2+z(1-e)+hez\},\\
 P_b={}&(1+z)(1+qz)+Sq(1+z)(2+z-2e)\\
       &+Sq^2z(1+z-e^2)+S^2q^2G^2.
\end{align*}
Then the exact rational identities are
\[
\begin{aligned}
 (1-h)\delta c_b&=-\frac{n(1-e)P_C}{zP},&
 h\delta a_b&=-\frac{n(1-e)P_A}{zP},\\
 \delta D_b&=-\frac{n(1-e)(1+qz)P_D}{zJP},&
 \delta h_2&=\frac{n(1-e)(1+qz)P_m}{zP},\\
 \delta h_b&=\frac{n(1-e)P_b}{zP},&
 \delta\alpha_b&=
 \frac{n(1-e)q^2(1+qz)[(1+S)e-S(1-h)(1+z)]}{zQP}.
\end{aligned}                                                   \tag{TU9}
\]
These identities are proved simply by substituting the six printed
RER1 weights and the six NC4 weights, clearing the indicated pure
source denominators, and multiplying.  They include $n=0$; the
right sides define no quotient by $n$.

After substituting $e=zw$, the polynomial
\[
       P_\Delta=\frac{P_b-(1+qz)P_m}{z}                 \tag{TU10}
\]
is a genuine finite polynomial in $S,z,q,w,h$.  Indeed the numerator
specialized at $z=0$ is zero identically, so ordinary monic division
by $z$ has zero remainder.  Consequently
\[
 \delta(h_b-h_2)=n(1-e)P_\Delta/P.
                                                               \tag{TU11}
\]
Every polynomial in (TU6), (TU9), (TU10) has bounded coefficients
and bounded arguments.  From (TU4) and (TU9)--(TU11) we obtain
\[
\begin{gathered}
 (1-h)\delta c_b,\ h\delta a_b,\ \delta D_b
                     =q\lambda_n\,O(1),\qquad
 \delta h_2,\delta h_b=q\lambda_n\,O(1),\\
 \delta\alpha_b/q^2=\lambda_n q\,O(S+e),\qquad
 \delta(h_b-h_2)/(qz)=\lambda_n O(1).
\end{gathered}                                                   \tag{TU12}
\]
Here $O(1)$ denotes an explicit bounded rational source array with
denominators among the units $P,J,Q,R$; it does not hide a state
or testing factor.  These conclusions require no lower bound on
$h$, $1-h$, or $M$.

In particular $c_b\asymp1/(1-h)$, $a_b\asymp1/h$,
$D_b\sim1/2$, $\alpha_b/q^2\sim1/2$, and
$h_2/q,h_b/q$ positive bounded units with common residue
$1/2+\overline\lambda_n$.  We also have
\[
 h_b/w=z/2,\qquad h_b-h_2=O(qz),\qquad
 H_0:=4h_2/R+4h_b/w\geq2z,
 \quad H_0\asymp q+z.                                  \tag{TU13}
\]
The last statements use the formal exact SC2 difference together
with (TU11), and hold also when $M$ is arbitrarily small.

\subsection{The complete additional current, including its normal}

Let the testing $u,v$ be positive infinitesimals and first take
$Z=1$ on its pure-testing state chart.  Use the exact SC coordinates
\[
 Q_b=(W+uY)/b,\quad Y_b=Y/S,\quad p=w/R,\quad
 N=((X+u)/b-pvQ_b)/g,
\]
\[
 H=uQ_b/b,\quad L=vQ_b/g,\quad
 \sigma=c_b+a_b,\quad \rho_0=Kh/R,\quad
 N_y=\sqrt\sigma\{D_uY_b-\rho_0H\}.
                                                               \tag{TU14}
\]
Both $N,N_y$ vanish at (TU7).  The exact SC3 full-parent identity
is an algebraic identity for arbitrary $c_b,a_b,D_b,\alpha_b,h_2,h_b$;
no C1 coefficient relation is used in deriving it.  With (TU8),
it represents all the actual C2 parent except its additional term
$4nD_uD_vX^2/S$.  Retain its original nonnegative $m$ packet
$4h_2wN^2/(MR)$, incorporate the second paired packet into
$H_0N^2$, and retain $4h_b(vN)^2/w$ exactly as in the uniform
QI normal lemma.  Define $\widehat N=\sqrt{H_0}N$.

Put $h_t=H/R$, and introduce the same four positive separated
monomials and their geometric mean as in QI:
\[
 A_*=ch_t,\quad B_*=cL/R,\quad C_*=uv/b,\quad
 D_*=qg,\qquad \Gamma=(A_*B_*C_*D_*)^{1/4}.
                                                               \tag{TU15}
\]
Their exact relations include
$\Gamma^2=qh_tv$ and $B_*C_*=(c/g)v^2h_t$.
The extra monomial
\[
                   U_*=cgu/b
\]
has the two exact identities
\[
       U_*B_*=\Gamma^2,\qquad
       cgX/b=czN+eB_*-U_*.
                                                               \tag{TU16}
\]
The entire additional current therefore has the following square
completion, with no state normal set to zero:
\begin{align*}
 \frac{4n}{S}D_uD_vX^2
 ={}&4\lambda_nD_uD_v(U_*-B_*-czN)^2\\
 &+4\lambda_nD_uD_v
   \{2(1-e)\Gamma^2-(1-e^2)B_*^2
                         -2(1-e)czB_*N\}.
                                                               \tag{TU17}
\end{align*}
The first line is a nonnegative separated SOS and vanishes at the
actual contact, including its unrestricted state-normal coordinate.
Its coefficient is literally zero if $n=0$.  The remaining scalar
changes have bounded coefficients on $B_*^2$ and $\Gamma^2$.
Their possibly nonzero real residues will cancel exactly with
the source comparison below.  The full additional first-normal coefficient, on
$B_*\widehat N$, is infinitesimal because
\[
       \lambda_n cz/\sqrt{H_0}=O(\lambda_n c\sqrt z)
                                            \in\mathrm I.
                                                               \tag{TU18}
\]
In particular (TU17) does not estimate or discard $u/b$, which can
be arbitrarily large.  The new positive square contains that entire
row, while the exact product relation in (TU16) places its cross
term on the old angular circuit.

\subsection{Why the full QI allocation is stable here}

We verify the actual finite hypotheses of QI1--QI10 rather than
an abstract continuity assertion.  Its coefficient arrays have
only the pure-testing denominators $D_u$ and positive pure-source
units after cancellation of the two $h$ openings.  The latter
cancellation is preserved exactly because
\[
 1-\rho_0=(1-h)/R,\qquad \rho_0=Kh/R.
                                                               \tag{TU19}
\]
For example all perturbations of
$c_b(1-\rho_0)^2$, $a_b\rho_0^2$,
$c_b(1-\rho_0)$ and $a_b\rho_0$ are $q\lambda_n$
times bounded rational arrays, by (TU9).  Moreover
\[
 \delta(\sigma\rho_0-c_b)
       =\rho_0\delta a_b-(1-\rho_0)\delta c_b
                             =q\lambda_n O(1).
                                                               \tag{TU20}
\]
This is the complete potentially singular harmonic first-normal
coefficient after multiplication by $\sqrt\sigma$.
Both $\sigma$ and its formal comparison are comparable to
$[h(1-h)]^{-1}$, and their ratio has residue one.
Their square-root ratios and inverses are consequently bounded
pure-source units.  Thus this change of hard normal loses no
opening bound.

In the exact generic scalar QI2, the changes of $a,b_*,c_*$ are
$q\lambda_n$ times bounded arrays in $u,v$.  This follows term by
term from (TU12), (TU19): its only $C,A$ occurrences are the first
two expressions in (TU19); its $D_b$ occurrences have bounded
factors; its $\alpha_b$ occurrences have bounded factors and the
additional small factors printed in QI2.  Consequently the
formal specialization $a_0,b_0,c_0$ and its exact $U,V$ divisions
in QI3--QI4 can be used unchanged, while the remaining $q$
coquotients change by bounded arrays.  Write
$\lambda_0=\overline\lambda_n$.  Their complete leading changes
are respectively $0,-8\lambda_0,4\lambda_0,0$: substituting
$S=z=q=u=v=0$ in the bounded coquotients (TU9) gives
$\delta a/q=-4\lambda_0$, $\delta b_*/q=0$,
$\delta c_*/q=-4\lambda_0$, and
$\delta(R^2P_L)/q=4\lambda_0$.  Thus before the extra current
the four residues are
\[
 \overline{a_1-c_1}=2,\qquad
 \overline{b_1+2c_1}=-8-8\lambda_0,\qquad
 \overline{R^2P_L/q}=2+4\lambda_0,\qquad
 \overline{4\alpha_bD_u/q^2}=2.
                                                               \tag{TU21}
\]
Equation (TU17) adds $-4\lambda_0$ to the $B_*^2$ residue
and $+8\lambda_0$ to the $\Gamma^2$ residue.  Consequently
the \emph{combined} full scalar has exactly the original four
residues $(2,-8,2,2)$ on $(A_*^2,\Gamma^2,B_*^2,D_*^2)$.
All differences from those residues are infinitesimal bounded
arrays, because $\lambda_n$ itself is bounded.  This cancellation
is why a thinness assumption is unnecessary on (TU1).
The one additional bound QI uses from physical $\beta\geq0$ is
also available directly, without any sign for $\beta^0$:
\[
  w/(c/g)=\sqrt{w/\tau}=O(1),\qquad w=O(1).
                                                               \tag{TU22}
\]
Thus the exact dangerous cross term $-4(wv/(c/g))D_uB_*C_*$
still has infinitesimal coefficient.

For the complete first $N$ normal, (TU13) gives exactly the two
bounds used in the uniform QI normal lemma:
$g/\sqrt{H_0}$ and $c/\sqrt{H_0}$ are bounded, and
\[
 \frac{h_bD_v-h_2}{\sqrt{qH_0}}=O(z+v^2).
                                                               \tag{TU23}
\]
The $H$ term and all other terms in SC14 therefore have the same
infinitesimal coefficients on the QI rows; (TU18) handles the whole
extra term.  For the harmonic normal, (TU20) adds a
$q\lambda_n$ bounded array to the left side of QI10.  It can be
absorbed into its existing $qK_q$ coquotient.  All other terms are
the same bounded rational arrays multiplied by the opening-safe
quantities in (TU12), (TU19).  The exact $D_b$ root and its two
normal terms in QI9 remain unchanged as algebraic identities.
The complete hard block therefore retains the real residue
$\operatorname{diag}(1,4)$ in $(\widehat N,N_y)$.

For clarity, here is the finite allocation being invoked.  Use
the QI roots
\[
 E_*=h_t-v,\quad F_*=vh_t-gD_*,\quad
 R_F=\frac{\sqrt{g\Gamma}}c(\Gamma-D_*),\quad
 R_E=\sqrt\Gamma(\sqrt{A_*}-\sqrt{cv}),
\]
together with the complete finite toric root list SC16.
The exact identities QI7--QI8 give strict real square pivots on
$F_*,R_F,R_E$; QI6 gives the strict $E_*$ pivot and twice the
four-monomial circuit.  Every remaining coefficient is an
infinitesimal bounded rational array in the testing $u,v$, with
the same finite denominators and divisions just verified.
All removed packets, including (TU17), vanish at (TU7).
The actual source contact value, its three scalar first jets,
and both full first-normal contact values consequently vanish
exactly for this remaining expression.  Applying the two finite
monic divisions SC17--SC20 and the displayed QI7 first-jet
identities therefore puts the entire scalar correction and
both complete first normals on this strict root list with
infinitesimal coefficients.  The coefficient functions are
independent of the state ratio $Q_b/Z$; no boundedness of that
ratio is assumed.  This is precisely the constructive QI
allocation, with its source inequalities checked in
(TU12)--(TU23), not a conclusion from a scalar restriction alone.

The resulting finite Gram is bounded and has positive definite
real residue.  The strict bounded-Gram lemma supplies tensor SOS.
Restore the original positive paired packets and the new full
square in (TU17).  This proves the local assertion for the entire
parent on the positive small-angle, negative-core chart.

\subsection{States, signs, and all remaining tests}

Every inverse and radical used above is pure source, pure testing,
or the explicitly separated positive monomial radical in (TU15).
For $Z\ne0$ the normalization $Z=1$ is pure testing; multiplying
back by $Z^2$ restores homogeneity.  At $Z=0$ the original
pure-testing core restriction is nonnegative.  The exact core
Schur identity and sign gauges used in the retained QI proof
force $W+uY>0$ on the remaining negative-core chart, and hence
all four monomials are positive.  Otherwise the core and all
original event packets already give the direct exit.  Literal
zero-angle charts have that same original exit.  Complex states
split into their real and imaginary parts, since the real
measurement matrices and full quadratic parent are unchanged.

The new R2 outer theorem R2O1--R2O24 treats every projective
testing pair outside the double-infinitesimal original chart,
with no condition on $a$ or $n/e^2$.  Its hypotheses are exactly
the original RER source conditions already retained here.
Combining that result with this local proof gives the claimed
all-tests, all-states theorem on (TU1).  No other source branch
is admitted by this conclusion.

\paragraph{Verification scope.}
The associated checker verifies the exact native contact map,
all six comparison formulas (TU9), the finite division (TU10),
the four complete residue cancellations in (TU21),
the whole generic SC reduction, the entire additional-current
completion (TU17), and its actual contact zero.  Ordered-field
bounds, the complete finite QI allocation and the outer-chart
join are analytical parts of this proof, not consequences of
the number of checker assertions.
